% 第九卷 9.1 节：同一输入“小王／阅读／通知”在五种结构预测范式下的输出表示（教学示意）。
% 行依次为序列标注（BMES）、跨度（嵌套区间）、转移系统（中间状态）、图（依存弧）、约束生成（括号序列前缀与允许集合）。
\begin{tikzpicture}[
  x=1mm,y=1mm,figure picture,
  rowlabel/.style={font=\figurefont\footnotesize,text=FigureInk,anchor=east,align=right},
  infer/.style={font=\figurefont\scriptsize,text=FigureMuted,anchor=west,align=left},
  ch/.style={rectangle,rounded corners=.6mm,draw=FigureInput,fill=FigureInputFill,
    line width=.6pt,minimum width=8mm,minimum height=6mm,inner sep=0pt,
    font=\figurefont\footnotesize},
  wd/.style={rectangle,rounded corners=.8mm,draw=FigureInput,fill=FigureInputFill,
    line width=.6pt,minimum width=16mm,minimum height=6mm,inner sep=0pt,
    font=\figurefont\footnotesize},
  tag/.style={font=\figurefont\footnotesize,text=BookTeal},
  span/.style={draw=FigureOutput,line width=.9pt},
  spanlab/.style={font=\figurefont\scriptsize,text=FigureInk,fill=white,inner sep=1pt},
  cfg/.style={rectangle,rounded corners=.8mm,draw=FigureModel,fill=white,
    line width=.6pt,minimum height=6.5mm,inner xsep=2pt,font=\figurefont\scriptsize},
  arc/.style={draw=BookTeal,line width=.9pt,-{Stealth[length=1.6mm,width=1.1mm]}},
  cand/.style={draw=FigureMuted,line width=.5pt,densely dashed,-{Stealth[length=1.4mm,width=1mm]}},
  sep/.style={draw=FigureGrid,line width=.4pt}
]
  % ---------- 行1：序列标注 ----------
  \node[rowlabel] at (20,0) {序列标注};
  \foreach \x/\c/\t in {30/小/B,39/王/E,55/阅/B,64/读/E,80/通/B,89/知/E}{
    \node[ch] at (\x,0) {\c};
    \node[tag] at (\x,-6.5) {\t};
  }
  \node[infer] at (122,-2) {Viterbi＋合法转移掩码\\$O(nL^2)$};
  \draw[sep] (0,-11) -- (158,-11);

  % ---------- 行2：跨度 ----------
  \node[rowlabel] at (20,-19) {跨度};
  \node[wd] at (34.5,-17) {小王};
  \node[wd] at (59.5,-17) {阅读};
  \node[wd] at (84.5,-17) {通知};
  \draw[span] (26.5,-22.5) -- ++(0,-1.2) -- (42.5,-23.7) -- ++(0,1.2);
  \node[spanlab,anchor=west] at (43.5,-23.4) {NP};
  \draw[span] (76.5,-22.5) -- ++(0,-1.2) -- (92.5,-23.7) -- ++(0,1.2);
  \node[spanlab,anchor=west] at (93.5,-23.4) {NP};
  \draw[span] (51.5,-25.5) -- ++(0,-1.4) -- (92.5,-26.9) -- ++(0,1.4);
  \node[spanlab,anchor=west] at (93.5,-26.7) {VP};
  \draw[span] (26.5,-28.6) -- ++(0,-1.6) -- (92.5,-30.2) -- ++(0,1.6);
  \node[spanlab,anchor=west] at (93.5,-30.0) {S};
  \node[infer] at (122,-21) {枚举候选$O(nW)$并剪枝\\不重叠选择或CKY};
  \draw[sep] (0,-34) -- (158,-34);

  % ---------- 行3：转移系统 ----------
  \node[rowlabel] at (20,-42) {转移系统};
  \node[cfg] (stk) at (37,-42) {栈：ROOT\enspace 阅读};
  \node[cfg] (buf) at (64,-42) {缓冲区：通知};
  \node[cfg] (arcs) at (91,-42) {已建：阅读$\to$小王};
  \node[font=\figurefont\scriptsize,text=BookTeal] at (64,-48.5) {下一动作：SHIFT};
  \node[infer] at (122,-42) {贪心或束搜索\\$2n$个动作};
  \draw[sep] (0,-53) -- (158,-53);

  % ---------- 行4：图与区间动态规划 ----------
  \node[rowlabel] at (20,-69) {图与区间DP};
  \node[wd,minimum width=12mm,draw=FigureMuted,fill=white] (r) at (30,-71) {ROOT};
  \node[wd,minimum width=13mm] (w1) at (50,-71) {小王};
  \node[wd,minimum width=13mm] (w2) at (71,-71) {阅读};
  \node[wd,minimum width=13mm] (w3) at (92,-71) {通知};
  \draw[arc] ([xshift=2mm]r.north) to[out=55,in=125] ([xshift=-1mm]w2.north);
  \draw[arc] ([xshift=-3mm]w2.north) to[out=120,in=60] ([xshift=2mm]w1.north);
  \draw[arc] ([xshift=3mm]w2.north) to[out=60,in=120] ([xshift=-2mm]w3.north);
  \draw[cand] (w1.south) to[out=-35,in=-145] (w3.south);
  \node[infer] at (122,-69) {Eisner $O(n^3)$\\Chu–Liu/Edmonds $O(n^2)$};
  \draw[sep] (0,-79) -- (158,-79);

  % ---------- 行5：约束生成 ----------
  \node[rowlabel] at (20,-87) {约束生成};
  \node[cfg,anchor=west,draw=FigureInput] (pre) at (24,-87)
    {已生成：(S (NP 小王) (VP 阅读 (NP 通知};
  \node[cfg,anchor=west,draw=BookTeal] (next) at ([xshift=4mm]pre.east)
    {下一词元只允许：)};
  \draw[figure flow] (pre.east) -- (next.west);
  \node[infer] at (122,-87) {束搜索＋文法掩码\\输出解析回结构};
\end{tikzpicture}
